\section{Layer 3: Bounding the ``Local Defect $\rightarrow$ Failure'' Story}
\label{sec:intervention}

A natural reading of the atlas is that a quantized agent makes a local error, and the task then fails unless the continuation recovers. We test this reading in two ways. First, a controlled intervention injects defects at frozen states. Second, a process audit inspects fork branches. Both results are \emph{bounding}: they limit what the simple story can explain, but they do not identify a mechanism.

\paragraph{Controlled defect injection.} We froze 30 ordinary states from q4 GPTQ-W4 trajectories, all from one development block. At each state we forced one of three actions as the next step: an \textsc{inv}alid action, a matched \textsc{val}id action, or a process-perturbation control (\textsc{pp}) that repeats the last successful read. Either B or Q then continued. The defect categories were wrong\_id, wrong\_tool and premature\_write, with 10 states each. % src: analysis/recovery/inject_q4/summary_v2.md; docs/RESEARCH_LOG.md W71
Every state passed a replay audit of 22 recorded fields (VAL$\times$Q against the original episode). Three fields were not recorded: the counters right after the restore, the rendered-prompt hash sequence and the tool-schema hash sequence. % src: analysis/recovery/inject_q4/summary_v2.md (Replay audit)

\begin{table}[t]
\centering\scriptsize\setlength{\tabcolsep}{2.8pt}
% src: analysis/recovery/inject_q4/summary_v2.md (Outcome turnover, population "all"; Controller dependence D)
\begin{tabular}{@{}llccl@{}}
\toprule
Ctrl & Contrast & resc./dam. & turnover & net [95\% CI] \\
\midrule
B & \textsc{inv} vs \textsc{val} & 4/3 & 7/30 & $+0.03$ [$-0.13$, $+0.20$] \\
B & \textsc{pp} vs \textsc{val} & 3/2 & 5/30 & $+0.03$ [$-0.10$, $+0.20$] \\
B & \textsc{inv} vs \textsc{pp} & 2/2 & 4/30 & $\phantom{+}0.00$ [$-0.13$, $+0.13$] \\
Q & \textsc{inv} vs \textsc{val} & 2/2 & 4/30 & $\phantom{+}0.00$ [$-0.13$, $+0.13$] \\
Q & \textsc{pp} vs \textsc{val} & 3/2 & 5/30 & $+0.03$ [$-0.10$, $+0.17$] \\
Q & \textsc{inv} vs \textsc{pp} & 2/3 & 5/30 & $-0.03$ [$-0.20$, $+0.10$] \\
\midrule
\multicolumn{4}{@{}l}{Controller dependence $D$ ($\neq0$ in 8/30 states)} & $+0.03$ [$-0.17$, $+0.27$] \\
\bottomrule
\end{tabular}
\caption{Defect injection on 30 frozen q4 states (development, one block): per-state rescue/damage of one arm against another, under the BF16 (B) or quantized (Q) continuation controller. $D=[Y(\textsc{inv},B)-Y(\textsc{val},B)]-[Y(\textsc{inv},Q)-Y(\textsc{val},Q)]$. Task-clustered bootstrap intervals; all exact McNemar $p\ge0.25$. Source: \nolinkurl{analysis/recovery/inject_q4/summary_v2.md}.}
\label{tab:recovery}
\end{table}

\paragraph{Result: two-sided turnover, no error-specific effect.} Injected defects change the outcome in 7/30 states under B (4 rescues, 3 damages) and 4/30 under Q (2 and 2). The process-perturbation control changes 5/30 under both controllers (Table~\ref{tab:recovery}). The error-specific effect, \textsc{inv} minus \textsc{pp}, is $0.00$ [$-0.13$, $+0.13$] under B and $-0.03$ [$-0.20$, $+0.10$] under Q. The controller dependence $D$ is nonzero in 8 of 30 states, but its mean is $+0.03$ [$-0.17$, $+0.27$]. % src: analysis/recovery/inject_q4/summary_v2.md
So injected defects cause heterogeneous two-sided turnover of the size a generic perturbation causes. We detect no error-specific effect beyond that perturbation and no mean controller difference, at $n=30$ states. An earlier reading of the same data, that injected defects ``seldom change outcomes'', was wrong and has been retracted.

\paragraph{Manipulation checks limit the categories.} Two of 30 states failed a manipulation check: an injected write was rejected, and a valid reference itself returned ``not found''. Restricting to the 28 passing states leaves the reading unchanged. The wrong\_tool category mostly produced argument errors: 8 of 10 injected calls failed with an unexpected-keyword error. We therefore call it a \emph{tool-call/schema error perturbation}, not a wrong-tool test. premature\_write is mostly failure$\rightarrow$failure: the valid arm under Q already fails in 7/10 of these states, and turnover is 0/10 under both controllers. % src: analysis/recovery/inject_q4/summary_v2.md (Manipulation checks; Outcome turnover by category); docs/RESEARCH_LOG.md W83
The event chain error $\rightarrow$ correction $\rightarrow$ related requirement satisfied completes in only 6 (B) and 7 (Q) of 29 injected errors, although 17 are corrected under each controller. Injections add tool errors ($+0.43$ under B, $+0.57$ under Q per branch). % src: analysis/recovery/inject_q4/summary_v2.md (Event chains; Paired differences)
\todo{RECOVERY-Q8: the same table for q8 GPTQ-W4 (30 states, block q3), from \nolinkurl{analysis/recovery/analyze_inject.py} on \nolinkurl{analysis/recovery/inject_q8}}

\paragraph{Process audit (LLM-annotated).} We sampled 80 development fork states, 20 per outcome sign. % src: manifests/annotation/wp5_sample_v1.json; analysis/annotation/wp5/table4_v2.md
 Two independent Sonnet 5.5 instances annotated them blind, a third instance adjudicated, and rule-based programmatic labels were added. These annotations are a \emph{process description} conditional on outcome-defined groups. They are not causal evidence, not prevalences and not human labels. Three patterns stand out (Appendix~\ref{app:table4}). First, local defects and final failure do not correspond one-to-one. In damage states, where Q's action led to success and B's to failure, Q's first-action defect is corrected later in 10 of 20 QQ branches. Second, whether a write executes does not separate the failing from the succeeding continuation of the same state (rescue 17 vs.\ 17; damage 15 vs.\ 14). Third, the failing continuation more often carries the rubric label \texttt{unintended\_or\_unconsented\_change} (rescue: QQ 9 vs.\ BQ 2, paired difference $+0.35$ [$+0.08$, $+0.57$]; damage: 1 vs.\ 9, $-0.40$ [$-0.62$, $-0.11$]). In damage states, 7 of these 20 BQ branches carry the label on writes \emph{with} recorded consent. That is the semantic ``wrong change'' part of the label, and it overlaps with the database-defined outcome by construction. Loops and budget exhaustion also mark failing branches (rescue: QQ 5 vs.\ BQ 0). % src: analysis/annotation/wp5/table4_v2.md
We make no claim about consent. The programmatic consent field agrees poorly with the annotators, so we use it only as a mechanical row. Re-labelling 30 states with fresh instances of the same model family reproduces the key pattern. That shows reproducibility, not correctness. % src: docs/RESEARCH_LOG.md W82 (analysis/annotation/wp5/human_check_v1/proxy_compare_v1.json)
